\chapter{Hướng dẫn cài đặt và chạy hệ thống}
\label{pl:caidat}

\ghichu{Hướng dẫn để người chấm hoặc người kế tiếp chạy được hệ thống trên máy của mình. Các bước dưới đây lấy từ \texttt{README.md} và \texttt{.env.example}; hãy chạy thử lại từ máy sạch và sửa nếu có bước nào không còn đúng. Không ghi giá trị bí mật thật vào báo cáo.}
\nguon{README.md, BE/.env.example}

\section{Yêu cầu}
Cần JDK 17 trở lên (đã chạy thử với JDK 21), Node.js và npm (đã chạy thử với Node.js 22.17), một dự án Supabase có PostgreSQL và Auth bật đăng nhập Google, và tùy chọn thông tin PayOS và khóa Google Gemini. Không cần cài Maven vì kho mã có sẵn Maven Wrapper.

\section{Chạy máy chủ}
Tạo tệp \texttt{BE/.env} từ mẫu \texttt{BE/.env.example}. Hai biến bắt buộc, thiếu là máy chủ không khởi động: \texttt{SUPABASE\_DB\_PASSWORD} và \texttt{APP\_JWT\_SECRET} (tối thiểu 48 byte, có thể sinh bằng \texttt{openssl rand -base64 48}). Sau đó chạy:

\begin{lstlisting}[language=bash]
cd BE
./mvnw spring-boot:run
\end{lstlisting}

Máy chủ lắng nghe ở \texttt{http://localhost:8080}. Khi khởi động, Flyway tự áp dụng các tệp trong \nolinkurl{BE/src/main/resources/db/migration}: cơ sở dữ liệu trống được tạo toàn bộ từ \texttt{V1\_\_baseline.sql}, cơ sở dữ liệu đã có dữ liệu được đánh dấu ở mốc \texttt{V1} và chỉ chạy các tệp mới hơn. Ở chế độ phát triển, \texttt{data.sql} nạp dữ liệu mẫu. Nếu không đặt thông tin PayOS, máy chủ chạy chế độ demo.

\section{Chạy giao diện}
Tạo tệp \texttt{FE/.env} với ba biến \texttt{VITE\_API\_BASE\_URL}, \texttt{VITE\_SUPABASE\_URL} và \texttt{VITE\_SUPABASE\_ANON\_KEY}, rồi chạy:

\begin{lstlisting}[language=bash]
cd FE
npm install
npm run dev
\end{lstlisting}

Giao diện chạy ở \texttt{http://localhost:5173}. Để dựng bản phát hành dùng \texttt{npm run build} (kiểm tra kiểu bằng \texttt{tsc} rồi đóng gói bằng Vite, kết quả ở thư mục \texttt{dist}).

\section{Chạy kiểm thử}
\begin{lstlisting}[language=bash]
cd BE
./mvnw test
\end{lstlisting}

Bộ kiểm thử không cần cơ sở dữ liệu; báo cáo chi tiết nằm ở \texttt{BE/target/surefire-reports}.

\chapter{Danh mục điểm cuối API}
\label{pl:api}

\ghichu{Bảng dưới được sinh tự động từ các lớp điều khiển bằng \texttt{report/DATN/tools/gen\_api\_appendix.py} (chạy lại lệnh này sau khi thêm hoặc đổi điểm cuối). Cột vai trò được phép gọi không sinh tự động vì phụ thuộc cả quy tắc theo đường dẫn ở \texttt{SecurityConfig} lẫn \texttt{@PreAuthorize}; xem Bảng \ref{tab:ht_route_quyen} ở Chương \ref{ch:hienthuc}. Nếu hội đồng cần, bổ sung cột này cho các nhóm chính.}

Bảng \ref{tab:pl_api} liệt kê toàn bộ điểm cuối của máy chủ.

% Sinh tu dong boi tools/gen_api_appendix.py - khong sua tay.
{\footnotesize
\begin{longtable}{|p{1.5cm}|p{7.3cm}|p{5.0cm}|}
\caption{Danh mục 133 điểm cuối API, nhóm theo lớp điều khiển}\label{tab:pl_api}\\
\hline
\textbf{Phương thức} & \textbf{Đường dẫn} & \textbf{Hàm xử lý} \\
\hline
\endfirsthead
\hline
\textbf{Phương thức} & \textbf{Đường dẫn} & \textbf{Hàm xử lý} \\
\hline
\endhead
\multicolumn{3}{|l|}{\textbf{\texttt{AdminController} (12)}} \\
\hline
GET & \texttt{/\allowbreak api/\allowbreak admin/\allowbreak branches} & \texttt{listBranches} \\
\hline
POST & \texttt{/\allowbreak api/\allowbreak admin/\allowbreak branches} & \texttt{createBranch} \\
\hline
PUT & \texttt{/\allowbreak api/\allowbreak admin/\allowbreak branches/\allowbreak \{id\}} & \texttt{updateBranch} \\
\hline
DELETE & \texttt{/\allowbreak api/\allowbreak admin/\allowbreak branches/\allowbreak \{id\}} & \texttt{deactivateBranch} \\
\hline
GET & \texttt{/\allowbreak api/\allowbreak admin/\allowbreak workspace-types} & \texttt{listWorkspaceTypes} \\
\hline
POST & \texttt{/\allowbreak api/\allowbreak admin/\allowbreak workspace-types} & \texttt{createWorkspaceType} \\
\hline
PUT & \texttt{/\allowbreak api/\allowbreak admin/\allowbreak workspace-types/\allowbreak \{id\}} & \texttt{updateWorkspaceType} \\
\hline
DELETE & \texttt{/\allowbreak api/\allowbreak admin/\allowbreak workspace-types/\allowbreak \{id\}} & \texttt{deleteWorkspaceType} \\
\hline
GET & \texttt{/\allowbreak api/\allowbreak admin/\allowbreak price-policies} & \texttt{listPricePolicies} \\
\hline
POST & \texttt{/\allowbreak api/\allowbreak admin/\allowbreak price-policies} & \texttt{createPricePolicy} \\
\hline
PUT & \texttt{/\allowbreak api/\allowbreak admin/\allowbreak price-policies/\allowbreak \{id\}} & \texttt{updatePricePolicy} \\
\hline
DELETE & \texttt{/\allowbreak api/\allowbreak admin/\allowbreak price-policies/\allowbreak \{id\}} & \texttt{deletePricePolicy} \\
\hline
\multicolumn{3}{|l|}{\textbf{\texttt{AdminLoyaltyController} (15)}} \\
\hline
GET & \texttt{/\allowbreak api/\allowbreak admin/\allowbreak amenities} & \texttt{listAmenities} \\
\hline
POST & \texttt{/\allowbreak api/\allowbreak admin/\allowbreak amenities} & \texttt{createAmenity} \\
\hline
PUT & \texttt{/\allowbreak api/\allowbreak admin/\allowbreak amenities/\allowbreak \{id\}} & \texttt{updateAmenity} \\
\hline
DELETE & \texttt{/\allowbreak api/\allowbreak admin/\allowbreak amenities/\allowbreak \{id\}} & \texttt{deleteAmenity} \\
\hline
GET & \texttt{/\allowbreak api/\allowbreak admin/\allowbreak workspace-type-amenities} & \texttt{listWorkspaceTypeAmenities} \\
\hline
PUT & \texttt{/\allowbreak api/\allowbreak admin/\allowbreak workspace-types/\allowbreak \{id\}/\allowbreak amenities} & \texttt{assignAmenities} \\
\hline
GET & \texttt{/\allowbreak api/\allowbreak admin/\allowbreak membership-tiers} & \texttt{listTiers} \\
\hline
POST & \texttt{/\allowbreak api/\allowbreak admin/\allowbreak membership-tiers} & \texttt{createTier} \\
\hline
PUT & \texttt{/\allowbreak api/\allowbreak admin/\allowbreak membership-tiers/\allowbreak \{id\}} & \texttt{updateTier} \\
\hline
DELETE & \texttt{/\allowbreak api/\allowbreak admin/\allowbreak membership-tiers/\allowbreak \{id\}} & \texttt{deleteTier} \\
\hline
POST & \texttt{/\allowbreak api/\allowbreak admin/\allowbreak membership-tiers/\allowbreak recalculate} & \texttt{recalculateTiers} \\
\hline
GET & \texttt{/\allowbreak api/\allowbreak admin/\allowbreak promotions} & \texttt{listPromotions} \\
\hline
POST & \texttt{/\allowbreak api/\allowbreak admin/\allowbreak promotions} & \texttt{createPromotion} \\
\hline
PUT & \texttt{/\allowbreak api/\allowbreak admin/\allowbreak promotions/\allowbreak \{id\}} & \texttt{updatePromotion} \\
\hline
DELETE & \texttt{/\allowbreak api/\allowbreak admin/\allowbreak promotions/\allowbreak \{id\}} & \texttt{deletePromotion} \\
\hline
\multicolumn{3}{|l|}{\textbf{\texttt{ApiHealthController} (1)}} \\
\hline
GET & \texttt{/\allowbreak api/\allowbreak health} & \texttt{health} \\
\hline
\multicolumn{3}{|l|}{\textbf{\texttt{AuditLogController} (1)}} \\
\hline
GET & \texttt{/\allowbreak api/\allowbreak admin/\allowbreak audit-logs} & \texttt{getAuditLogs} \\
\hline
\multicolumn{3}{|l|}{\textbf{\texttt{AuthController} (4)}} \\
\hline
POST & \texttt{/\allowbreak api/\allowbreak auth/\allowbreak refresh} & \texttt{refresh} \\
\hline
POST & \texttt{/\allowbreak api/\allowbreak auth/\allowbreak register} & \texttt{register} \\
\hline
POST & \texttt{/\allowbreak api/\allowbreak auth/\allowbreak login} & \texttt{login} \\
\hline
POST & \texttt{/\allowbreak api/\allowbreak auth/\allowbreak sync} & \texttt{sync} \\
\hline
\multicolumn{3}{|l|}{\textbf{\texttt{BookingAddonController} (4)}} \\
\hline
GET & \texttt{/\allowbreak api/\allowbreak bookings/\allowbreak \{bookingId\}/\allowbreak addons} & \texttt{getTab} \\
\hline
POST & \texttt{/\allowbreak api/\allowbreak bookings/\allowbreak \{bookingId\}/\allowbreak addons} & \texttt{addService} \\
\hline
DELETE & \texttt{/\allowbreak api/\allowbreak bookings/\allowbreak \{bookingId\}/\allowbreak addons/\allowbreak \{itemId\}} & \texttt{voidItem} \\
\hline
POST & \texttt{/\allowbreak api/\allowbreak bookings/\allowbreak \{bookingId\}/\allowbreak addons/\allowbreak settle} & \texttt{settle} \\
\hline
\multicolumn{3}{|l|}{\textbf{\texttt{BookingController} (6)}} \\
\hline
POST & \texttt{/\allowbreak api/\allowbreak bookings} & \texttt{create} \\
\hline
POST & \texttt{/\allowbreak api/\allowbreak bookings/\allowbreak quote} & \texttt{quote} \\
\hline
GET & \texttt{/\allowbreak api/\allowbreak bookings/\allowbreak my} & \texttt{myBookings} \\
\hline
GET & \texttt{/\allowbreak api/\allowbreak bookings/\allowbreak \{id\}} & \texttt{getOne} \\
\hline
GET & \texttt{/\allowbreak api/\allowbreak bookings/\allowbreak code/\allowbreak \{code\}} & \texttt{getByCode} \\
\hline
GET & \texttt{/\allowbreak api/\allowbreak bookings/\allowbreak branch-today} & \texttt{getBranchTodayBookings} \\
\hline
\multicolumn{3}{|l|}{\textbf{\texttt{BranchAdminSpaceController} (18)}} \\
\hline
GET & \texttt{/\allowbreak api/\allowbreak branch-admin/\allowbreak branches/\allowbreak \{branchId\}/\allowbreak name} & \texttt{getBranchName} \\
\hline
GET & \texttt{/\allowbreak api/\allowbreak branch-admin/\allowbreak workspace-types} & \texttt{listWorkspaceTypes} \\
\hline
GET & \texttt{/\allowbreak api/\allowbreak branch-admin/\allowbreak floors} & \texttt{listFloors} \\
\hline
POST & \texttt{/\allowbreak api/\allowbreak branch-admin/\allowbreak floors} & \texttt{createFloor} \\
\hline
PUT & \texttt{/\allowbreak api/\allowbreak branch-admin/\allowbreak floors/\allowbreak \{id\}} & \texttt{updateFloor} \\
\hline
DELETE & \texttt{/\allowbreak api/\allowbreak branch-admin/\allowbreak floors/\allowbreak \{id\}} & \texttt{deleteFloor} \\
\hline
GET & \texttt{/\allowbreak api/\allowbreak branch-admin/\allowbreak floors/\allowbreak \{floorId\}/\allowbreak workspaces} & \texttt{listWorkspaces} \\
\hline
POST & \texttt{/\allowbreak api/\allowbreak branch-admin/\allowbreak workspaces} & \texttt{createWorkspace} \\
\hline
PUT & \texttt{/\allowbreak api/\allowbreak branch-admin/\allowbreak workspaces/\allowbreak \{id\}} & \texttt{updateWorkspace} \\
\hline
DELETE & \texttt{/\allowbreak api/\allowbreak branch-admin/\allowbreak workspaces/\allowbreak \{id\}} & \texttt{deleteWorkspace} \\
\hline
GET & \texttt{/\allowbreak api/\allowbreak branch-admin/\allowbreak staff} & \texttt{listStaff} \\
\hline
POST & \texttt{/\allowbreak api/\allowbreak branch-admin/\allowbreak staff} & \texttt{createStaff} \\
\hline
PUT & \texttt{/\allowbreak api/\allowbreak branch-admin/\allowbreak staff/\allowbreak \{staffId\}} & \texttt{updateStaff} \\
\hline
PUT & \texttt{/\allowbreak api/\allowbreak branch-admin/\allowbreak staff/\allowbreak \{staffId\}/\allowbreak status} & \texttt{updateStaffStatus} \\
\hline
GET & \texttt{/\allowbreak api/\allowbreak branch-admin/\allowbreak price-policies} & \texttt{listPricePolicies} \\
\hline
POST & \texttt{/\allowbreak api/\allowbreak branch-admin/\allowbreak price-policies} & \texttt{createPricePolicy} \\
\hline
PUT & \texttt{/\allowbreak api/\allowbreak branch-admin/\allowbreak price-policies/\allowbreak \{id\}} & \texttt{updatePricePolicy} \\
\hline
DELETE & \texttt{/\allowbreak api/\allowbreak branch-admin/\allowbreak price-policies/\allowbreak \{id\}} & \texttt{deletePricePolicy} \\
\hline
\multicolumn{3}{|l|}{\textbf{\texttt{CancellationPolicyController} (5)}} \\
\hline
GET & \texttt{/\allowbreak api/\allowbreak cancellation-policies} & \texttt{getPolicies} \\
\hline
POST & \texttt{/\allowbreak api/\allowbreak cancellation-policies} & \texttt{createPolicy} \\
\hline
PUT & \texttt{/\allowbreak api/\allowbreak cancellation-policies/\allowbreak \{id\}} & \texttt{updatePolicy} \\
\hline
DELETE & \texttt{/\allowbreak api/\allowbreak cancellation-policies/\allowbreak \{id\}} & \texttt{deletePolicy} \\
\hline
POST & \texttt{/\allowbreak api/\allowbreak bookings/\allowbreak \{bookingId\}/\allowbreak cancel-v2} & \texttt{cancelBookingWithRefund} \\
\hline
\multicolumn{3}{|l|}{\textbf{\texttt{ChatbotController} (2)}} \\
\hline
POST & \texttt{/\allowbreak api/\allowbreak chatbot/\allowbreak message} & \texttt{message} \\
\hline
POST & \texttt{/\allowbreak api/\allowbreak chatbot/\allowbreak actions/\allowbreak confirm} & \texttt{confirm} \\
\hline
\multicolumn{3}{|l|}{\textbf{\texttt{CheckinController} (3)}} \\
\hline
POST & \texttt{/\allowbreak api/\allowbreak checkins/\allowbreak booking/\allowbreak \{bookingId\}} & \texttt{checkin} \\
\hline
POST & \texttt{/\allowbreak api/\allowbreak checkins/\allowbreak \{id\}/\allowbreak checkout} & \texttt{checkout} \\
\hline
GET & \texttt{/\allowbreak api/\allowbreak checkins/\allowbreak active} & \texttt{getActiveCheckins} \\
\hline
\multicolumn{3}{|l|}{\textbf{\texttt{CommunityPostController} (4)}} \\
\hline
GET & \texttt{/\allowbreak api/\allowbreak community/\allowbreak posts} & \texttt{listFeed} \\
\hline
POST & \texttt{/\allowbreak api/\allowbreak community/\allowbreak posts} & \texttt{createPost} \\
\hline
DELETE & \texttt{/\allowbreak api/\allowbreak community/\allowbreak posts/\allowbreak \{postId\}} & \texttt{deletePost} \\
\hline
GET & \texttt{/\allowbreak api/\allowbreak community/\allowbreak tags} & \texttt{listTags} \\
\hline
\multicolumn{3}{|l|}{\textbf{\texttt{CustomerSpaceController} (6)}} \\
\hline
GET & \texttt{/\allowbreak api/\allowbreak customer/\allowbreak spaces/\allowbreak branches} & \texttt{listActiveBranches} \\
\hline
GET & \texttt{/\allowbreak api/\allowbreak customer/\allowbreak spaces/\allowbreak branches/\allowbreak \{branchId\}/\allowbreak prices} & \texttt{listPrices} \\
\hline
GET & \texttt{/\allowbreak api/\allowbreak customer/\allowbreak spaces/\allowbreak pricing-summary} & \texttt{pricingSummary} \\
\hline
GET & \texttt{/\allowbreak api/\allowbreak customer/\allowbreak spaces/\allowbreak branches/\allowbreak \{branchId\}/\allowbreak floors} & \texttt{listFloors} \\
\hline
GET & \texttt{/\allowbreak api/\allowbreak customer/\allowbreak spaces/\allowbreak branches/\allowbreak \{branchId\}/\allowbreak floors/\allowbreak \{floorId\}/\allowbreak workspaces} & \texttt{listWorkspaces} \\
\hline
GET & \texttt{/\allowbreak api/\allowbreak customer/\allowbreak spaces/\allowbreak branches/\allowbreak \{branchId\}/\allowbreak booking-status} & \texttt{getBookingStatus} \\
\hline
\multicolumn{3}{|l|}{\textbf{\texttt{ExtraServiceController} (5)}} \\
\hline
GET & \texttt{/\allowbreak api/\allowbreak extra-services} & \texttt{getAvailableServices} \\
\hline
GET & \texttt{/\allowbreak api/\allowbreak extra-services/\allowbreak all} & \texttt{getAllServices} \\
\hline
POST & \texttt{/\allowbreak api/\allowbreak extra-services} & \texttt{createService} \\
\hline
PUT & \texttt{/\allowbreak api/\allowbreak extra-services/\allowbreak \{id\}} & \texttt{updateService} \\
\hline
DELETE & \texttt{/\allowbreak api/\allowbreak extra-services/\allowbreak \{id\}} & \texttt{deleteService} \\
\hline
\multicolumn{3}{|l|}{\textbf{\texttt{LoyaltyController} (4)}} \\
\hline
GET & \texttt{/\allowbreak api/\allowbreak customer/\allowbreak spaces/\allowbreak workspace-types} & \texttt{listWorkspaceTypesWithAmenities} \\
\hline
GET & \texttt{/\allowbreak api/\allowbreak membership/\allowbreak tiers} & \texttt{listTiers} \\
\hline
GET & \texttt{/\allowbreak api/\allowbreak membership/\allowbreak me} & \texttt{myMembership} \\
\hline
GET & \texttt{/\allowbreak api/\allowbreak promotions/\allowbreak available} & \texttt{availablePromotions} \\
\hline
\multicolumn{3}{|l|}{\textbf{\texttt{NotificationController} (4)}} \\
\hline
GET & \texttt{/\allowbreak api/\allowbreak notifications} & \texttt{getMyNotifications} \\
\hline
GET & \texttt{/\allowbreak api/\allowbreak notifications/\allowbreak unread-count} & \texttt{getUnreadCount} \\
\hline
PATCH & \texttt{/\allowbreak api/\allowbreak notifications/\allowbreak \{id\}/\allowbreak read} & \texttt{markAsRead} \\
\hline
PATCH & \texttt{/\allowbreak api/\allowbreak notifications/\allowbreak read-all} & \texttt{markAllAsRead} \\
\hline
\multicolumn{3}{|l|}{\textbf{\texttt{PartnerMatchingController} (3)}} \\
\hline
GET & \texttt{/\allowbreak api/\allowbreak matching/\allowbreak suggestions} & \texttt{getSuggestedPartners} \\
\hline
GET & \texttt{/\allowbreak api/\allowbreak profiles/\allowbreak me/\allowbreak networking} & \texttt{getMyNetworkingProfile} \\
\hline
PUT & \texttt{/\allowbreak api/\allowbreak profiles/\allowbreak me/\allowbreak networking} & \texttt{updateMyNetworkingProfile} \\
\hline
\multicolumn{3}{|l|}{\textbf{\texttt{PaymentController} (9)}} \\
\hline
POST & \texttt{/\allowbreak api/\allowbreak payments/\allowbreak momo/\allowbreak create} & \texttt{createMomoPayment} \\
\hline
POST & \texttt{/\allowbreak api/\allowbreak payments/\allowbreak payos/\allowbreak create} & \texttt{createPayosPayment} \\
\hline
POST & \texttt{/\allowbreak api/\allowbreak payments/\allowbreak cash/\allowbreak create} & \texttt{createCashPayment} \\
\hline
GET & \texttt{/\allowbreak api/\allowbreak payments/\allowbreak booking/\allowbreak \{bookingId\}} & \texttt{paymentsForBooking} \\
\hline
POST & \texttt{/\allowbreak api/\allowbreak payments/\allowbreak momo/\allowbreak notify} & \texttt{momoIpn} \\
\hline
GET & \texttt{/\allowbreak api/\allowbreak payments/\allowbreak momo/\allowbreak return} & \texttt{momoReturn} \\
\hline
POST & \texttt{/\allowbreak api/\allowbreak payments/\allowbreak payos/\allowbreak notify} & \texttt{payosWebhook} \\
\hline
GET & \texttt{/\allowbreak api/\allowbreak payments/\allowbreak payos/\allowbreak return} & \texttt{payosReturn} \\
\hline
GET & \texttt{/\allowbreak api/\allowbreak payments/\allowbreak payos/\allowbreak status/\allowbreak \{orderCode\}} & \texttt{getPayosPaymentStatus} \\
\hline
\multicolumn{3}{|l|}{\textbf{\texttt{PayosSimulationController} (1)}} \\
\hline
POST & \texttt{/\allowbreak api/\allowbreak payments/\allowbreak payos/\allowbreak simulate} & \texttt{simulatePayosPayment} \\
\hline
\multicolumn{3}{|l|}{\textbf{\texttt{RefundController} (3)}} \\
\hline
GET & \texttt{/\allowbreak api/\allowbreak refunds} & \texttt{list} \\
\hline
POST & \texttt{/\allowbreak api/\allowbreak refunds/\allowbreak \{id\}/\allowbreak process} & \texttt{process} \\
\hline
POST & \texttt{/\allowbreak api/\allowbreak refunds/\allowbreak \{id\}/\allowbreak reject} & \texttt{reject} \\
\hline
\multicolumn{3}{|l|}{\textbf{\texttt{ReportController} (2)}} \\
\hline
GET & \texttt{/\allowbreak api/\allowbreak reports/\allowbreak overview} & \texttt{getOverview} \\
\hline
GET & \texttt{/\allowbreak api/\allowbreak reports/\allowbreak export/\allowbreak csv} & \texttt{exportCsv} \\
\hline
\multicolumn{3}{|l|}{\textbf{\texttt{StaffBookingController} (3)}} \\
\hline
POST & \texttt{/\allowbreak api/\allowbreak staff/\allowbreak bookings/\allowbreak \{bookingId\}/\allowbreak cancel} & \texttt{cancelForCustomer} \\
\hline
POST & \texttt{/\allowbreak api/\allowbreak staff/\allowbreak bookings/\allowbreak walkin} & \texttt{createWalkinBooking} \\
\hline
GET & \texttt{/\allowbreak api/\allowbreak staff/\allowbreak bookings/\allowbreak branches/\allowbreak \{branchId\}/\allowbreak workspaces-booking-status} & \texttt{getWorkspaceBookingStatus} \\
\hline
\multicolumn{3}{|l|}{\textbf{\texttt{StaffDashboardController} (1)}} \\
\hline
GET & \texttt{/\allowbreak api/\allowbreak staff/\allowbreak dashboard/\allowbreak stats} & \texttt{getStats} \\
\hline
\multicolumn{3}{|l|}{\textbf{\texttt{StaffMaintenanceController} (5)}} \\
\hline
GET & \texttt{/\allowbreak api/\allowbreak staff/\allowbreak branches/\allowbreak \{branchId\}/\allowbreak maintenance} & \texttt{getAllMaintenanceForBranch} \\
\hline
GET & \texttt{/\allowbreak api/\allowbreak staff/\allowbreak branches/\allowbreak \{branchId\}/\allowbreak workspaces-maintenance} & \texttt{getWorkspaceMaintenanceStatus} \\
\hline
POST & \texttt{/\allowbreak api/\allowbreak staff/\allowbreak workspaces/\allowbreak \{workspaceId\}/\allowbreak maintenance} & \texttt{createMaintenance} \\
\hline
PUT & \texttt{/\allowbreak api/\allowbreak staff/\allowbreak maintenance/\allowbreak \{maintenanceId\}/\allowbreak complete} & \texttt{completeMaintenance} \\
\hline
DELETE & \texttt{/\allowbreak api/\allowbreak staff/\allowbreak maintenance/\allowbreak \{maintenanceId\}} & \texttt{deleteMaintenance} \\
\hline
\multicolumn{3}{|l|}{\textbf{\texttt{TagController} (4)}} \\
\hline
GET & \texttt{/\allowbreak api/\allowbreak tags} & \texttt{getTags} \\
\hline
POST & \texttt{/\allowbreak api/\allowbreak tags} & \texttt{createTag} \\
\hline
PUT & \texttt{/\allowbreak api/\allowbreak tags/\allowbreak \{id\}} & \texttt{updateTag} \\
\hline
DELETE & \texttt{/\allowbreak api/\allowbreak tags/\allowbreak \{id\}} & \texttt{deleteTag} \\
\hline
\multicolumn{3}{|l|}{\textbf{\texttt{UserController} (8)}} \\
\hline
GET & \texttt{/\allowbreak api/\allowbreak users/\allowbreak profile} & \texttt{getProfile} \\
\hline
PUT & \texttt{/\allowbreak api/\allowbreak users/\allowbreak profile} & \texttt{updateProfile} \\
\hline
PUT & \texttt{/\allowbreak api/\allowbreak users/\allowbreak change-password} & \texttt{changePassword} \\
\hline
GET & \texttt{/\allowbreak api/\allowbreak users/\allowbreak search} & \texttt{searchUsers} \\
\hline
GET & \texttt{/\allowbreak api/\allowbreak users} & \texttt{getUsers} \\
\hline
PUT & \texttt{/\allowbreak api/\allowbreak users/\allowbreak \{id\}/\allowbreak status} & \texttt{updateUserStatus} \\
\hline
PUT & \texttt{/\allowbreak api/\allowbreak users/\allowbreak \{id\}/\allowbreak role} & \texttt{updateUserRole} \\
\hline
POST & \texttt{/\allowbreak api/\allowbreak users/\allowbreak walkin} & \texttt{createWalkinUser} \\
\hline
\end{longtable}
}


\chapter{Các quyết định nghiệp vụ đã chốt sau đợt rà soát}
\label{pl:quyetdinh}

\ghichu{Bảng này tóm tắt mục 5.1 của bản rà soát ngày 18/09/2026, các quyết định được chủ dự án chốt ngày 20/09/2026. Giữ nguyên mã Q để đối chiếu với tài liệu gốc. Có thể dẫn bảng này từ mục quy tắc nghiệp vụ ở Chương \ref{ch:yeucau}. Lưu ý Q4: bản rà soát ghi hàng chờ hoàn tiền do ``staff hoặc branch admin'' duyệt, nhưng \texttt{RefundController} chỉ cho quản lý chi nhánh và quản trị hệ thống (\texttt{@PreAuthorize}); hãy thống nhất lại câu chữ.}
\nguon{report/LOGIC\_AUDIT.md (muc 5.1)}

\begin{table}[H]
\centering
\caption{Các quyết định nghiệp vụ đã chốt ngày 20/09/2026}
\label{tab:pl_quyetdinh}
\footnotesize
\renewcommand{\arraystretch}{1.2}
\begin{tabular}{|p{1.1cm}|p{7.6cm}|p{5.1cm}|}
\hline
\textbf{Mã} & \textbf{Quyết định} & \textbf{Ảnh hưởng} \\
\hline
Q1 & Chặn hủy sau giờ bắt đầu; khách không đến thành không đến (\texttt{NO\_SHOW}) và mất phí & Xử lý cùng lúc ba lỗi hoàn tiền sai của mã \texttt{CAN-01} \\
\hline
Q2 & Khoảng thời gian của chính sách hủy tính theo phút, nửa khoảng \texttt{[min, max)} & Hết chồng biên tại đúng 72 giờ (\texttt{CAN-04}) \\
\hline
Q3 & Không khớp chính sách chi nhánh thì rơi về chính sách toàn cục & Không sửa mã; ghi rõ vào tài liệu \\
\hline
Q4 & Giữ hàng chờ hoàn tiền do người có quyền duyệt xử lý, hệ thống không tự chi trả & Không sửa mã; sửa tài liệu \\
\hline
Q5 & Bỏ tùy chọn tiền mặt ở giao diện khách hàng; chỉ nhân viên tạo đơn thu tiền mặt & Sửa trang xác nhận đặt chỗ (\texttt{PAY-03}) \\
\hline
Q5b & Đơn tạo tại quầy giữ hạn thanh toán 15 phút như đơn trực tuyến & Không sửa mã (\texttt{BKG-03}) \\
\hline
Q6 & Bảo trì chỉ hoàn phần thời gian giao nhau, không hủy cả hợp đồng; chặn bảo trì bắt đầu ở quá khứ & Sửa \texttt{StaffMaintenanceService} (\texttt{MT-01}) \\
\hline
Q7 & Giữ trạng thái \texttt{NO\_SHOW} riêng, không gộp vào \texttt{COMPLETED} & Không sửa mã; sửa tài liệu \\
\hline
Q8 & Đơn đang sử dụng (\texttt{CHECKED\_IN}) không hủy được, khách muốn về sớm thì check-out & Không sửa mã; sửa tài liệu \\
\hline
Q9 & Đơn hết hạn thì trả lại lượt mã khuyến mãi, đơn đã thanh toán rồi hủy thì không & Sửa truy vấn đếm lượt dùng mã (\texttt{PRM-01}) \\
\hline
Q10 & Hạng thành viên tính theo chi tiêu của đơn đã dùng, trừ phần đã hoàn & Sửa \texttt{MembershipService} (\texttt{MEM-01}) \\
\hline
Q11 & Doanh thu là tổng tiền các đơn đã dùng (\texttt{COMPLETED}, \texttt{NO\_SHOW}) nhóm theo giờ kết thúc, trừ khoản đã hoàn và dịch vụ chưa thu & Sửa \texttt{ReportService} (\texttt{RPT-01}) \\
\hline
Q12 & Giữ công thức gợi ý đối tác trong mã (0,5; 0,2; 0,2; cộng 0,15 cùng chi nhánh), bỏ sàn hiển thị, chỉ gợi ý khách hàng & Sửa mã hai chỗ; sửa tài liệu (\texttt{MAT-01}, \texttt{MAT-03}) \\
\hline
Q13 & Nhân viên được hủy thay khách: bắt buộc lý do, có nhật ký kiểm toán, được miễn phạt cho ca ngoại lệ & Thêm điểm cuối hủy thay khách \\
\hline
Q14 & Cửa sổ check-in kết thúc ở giờ kết thúc của đơn & Không sửa mã; sửa tài liệu (\texttt{CHK-02}) \\
\hline
Q15 & Lấy lược đồ đang chạy làm nền, quản lý bằng Flyway; các tệp di trú cũ chuyển vào kho lưu trữ & Thêm Flyway, \texttt{V1} và \texttt{V2} (\texttt{X-02}) \\
\hline
Q16 & Giới hạn thời lượng đặt chỗ: 24 giờ, 30 ngày, 52 tuần, 12 tháng & Sửa \texttt{computeUnitCount} (\texttt{BKG-02}) \\
\hline
\end{tabular}
\end{table}
